\originalchapter{8.324 作业7题解}
本章根据 MIT OCW 8.324(2010年秋)Assignment 7 与官方解答中文重编并补充计算.原题1与第6套原题4重复,仍保留本套题号;原题3仅为商书题号.采用本书 $(+,-,-,-)$;原资料的六维标量作用量使用相反度规,转写后为正时向动能,$\phi^3$ 相互作用的符号按题单保留.

\originalsection{原题1}
\begin{exercise}
重整化群性质,原题30分.(a) 对可逆耦合变换 $\bar\lambda=f(\lambda)$ 求 $\bar\beta$,说明张量类型;(b) 对 $\beta=b_2\lambda^2+b_3\lambda^3+b_4\lambda^4+\cdots$、$f(\lambda)=\lambda+a_2\lambda^2+a_3\lambda^3+\cdots$,求新系数,讨论前两系数不变、$\bar b_4$ 的任意调整,以及固定点和其导数;(c) 对 $\mu\,\dd g/\dd\mu=-bg^2-cg^3-dg^4-\cdots$ 求弱耦合积分式及 RG 不变尺度;(d) 证明无其他尺度的单耦合理论的维度转化.
\end{exercise}
\begin{solution}
(a) $\bar\beta=\dd\bar\lambda/\dd\log\mu=f'(\lambda)\beta$,是耦合空间上的逆变切向量.

(b) 为避免只给方案变化公式,先列出比较所需的两组展开:
\begin{align*}
f'\beta={}&b_2\lambda^2+(b_3+2a_2b_2)\lambda^3
+(b_4+2a_2b_3+3a_3b_2)\lambda^4+\cdots,\\
\bar\lambda^2={}&\lambda^2+2a_2\lambda^3+(a_2^2+2a_3)\lambda^4+\cdots,\\
\bar\lambda^3={}&\lambda^3+3a_2\lambda^4+\cdots.
\end{align*}
逐次匹配 $\bar\beta=\bar b_2\bar\lambda^2+\bar b_3\bar\lambda^3+\bar b_4\bar\lambda^4+\cdots$,得
\[
\bar b_2=b_2,\qquad\bar b_3=b_3,
\qquad\bar b_4=b_4-a_2b_3+(a_3-a_2^2)b_2.
\]
若 $b_2\ne0$ 可选 $a_3$;若 $b_2=0,b_3\ne0$ 可选 $a_2$;若两者都为零则首个非零 $b_4$ 不能任意改写.固定点随可逆坐标变换映到固定点,因为 $f'\beta=0$;求导得
\[
\bar\beta'=\frac{\dd(f'\beta)}{\dd\bar\lambda}
=\beta'+\frac{f''}{f'}\beta,
\qquad\bar\beta'(f(\lambda_F))=\beta'(\lambda_F).
\]
这里要求 $f'(\lambda_F)$ 有限非零;奇异的“坐标变换”不适用.

(c) 令 $t=\log\mu$,倒置微分方程得
\[
\dd t=\left[-\frac1{bg^2}+\frac c{b^2g}+\order(1)\right]\dd g,
\qquad
\log\frac\mu\Lambda=\frac1{bg(\mu)}+
\frac c{b^2}\log(bg(\mu))+\order(g(\mu)).
\]
取 $b\ne0,bg>0$;若是其他实分支用 $\log|bg|$.以 $F'(g)=1/\beta(g)$ 精确定义 $\Lambda=\mu\e^{-F(g)}$,其对数导数为 $1-F'\beta=0$.故 $\Lambda$ 是流的积分常数,弱耦合式给出它的渐近表达.

(d) 没有其他尺度、且只有一个耦合时,$\beta$ 是 $g$ 的函数;非固定点分支上积分给 $F(g)=\log(\mu/\Lambda)$,局部反演给 $g=f_0(\log(\mu/\Lambda))$.不同理论有不同的 $F$ 及 $f_0$.无量纲边界条件转成有量纲积分常数,这就是维度转化.固定点分支是常数耦合,须与这一可逆积分分支区分.
\end{solution}

\originalsection{原题2}
\begin{exercise}
六维实标量理论的裸作用量为
\[
\Lag_0=\frac12(\partial\phi_0)^2-\frac12m_0^2\phi_0^2-\frac{g_0}{6}\phi_0^3.
\]
用维数正规化与最小减除求 $g$ 的 $\beta$ 函数,说明高能耦合减弱.原题30分.
\end{exercise}
\begin{solution}
取 $d=6-\epsilon$,$\phi_0=Z_\phi^{1/2}\phi$,$Z_\phi=1+A$,将顶点反项写为 $g\mu^{\epsilon/2}(1+C)$,则
\[
g_0=\mu^{\epsilon/2}g(1+C)(1+A)^{-3/2}.
\]
三个外腿的场重整化因子必须保留;只求三角图而漏掉它会得到错误的 $\beta$ 系数.定义 $\alpha=g^2/(4\pi)^3$.

一圈两点图的极点,可由 $1/(ab)$ 参数化和 Wick 转动得到.用源资料的空间样变量 $p_{\rm source}^2=-p^2$ 写 $\Delta=m^2+x(1-x)p_{\rm source}^2$,其极点部分为
\[
\frac{g^2\mu^\epsilon}{2(4\pi)^{d/2}}
\Gamma(2-d/2)\int_0^1\Delta^{d/2-2}\dd x
=-\frac\alpha\epsilon\left(m^2+\frac{p_{\rm source}^2}{6}\right)+\text{有限项}.
\]
因为 $\Gamma(-1+\epsilon/2)=-2/\epsilon+\order(1)$,$\int_0^1x(1-x)\dd x=1/6$,动量相关反项遂为
\[
A=-\frac\alpha{6\epsilon}+\order(g^4).
\]
质量反项影响 $m_0$,但不进入本阶的 $g_0$.一圈三点图与第4套原题2相同,其参数三角形面积 $1/2$、Gamma 极点 $2/\epsilon$,因此顶点发散相对于树顶点为 $\alpha/\epsilon$.MS 只消去此极点,取
\[
C=-\frac\alpha\epsilon+\order(g^4).
\]
相互作用符号整体改变等价于 $g\mapsto-g$,相对的一圈系数依赖 $g^2$,不改变以下结论.

合并并展开到三阶,
\[
g_0=\mu^{\epsilon/2}g
\left(1-\frac\alpha\epsilon\right)
\left(1-\frac\alpha{6\epsilon}\right)^{-3/2}
=\mu^{\epsilon/2}\left[g-
\frac{3g^3}{4(4\pi)^3\epsilon}+\order(g^5)\right].
\]
记 $c_0=3/[4(4\pi)^3]$,裸耦合与 $\mu$ 无关,故
\[
0=\frac\epsilon2\left(g-\frac{c_0g^3}{\epsilon}\right)
+\beta_g\left(1-\frac{3c_0g^2}{\epsilon}\right).
\]
先保留树级 $\beta_g=-\epsilon g/2+\beta_3$,再代入比较三阶项:$0=c_0g^3+\beta_3$.因此
\[
\beta_g=-\frac\epsilon2g-\frac{3g^3}{4(4\pi)^3}+\order(g^5),
\qquad
\beta_\alpha=-\epsilon\alpha-\frac32\alpha^2+\order(\alpha^3).
\]
在真正六维 $\epsilon=0$ 时,积分一圈方程得
\[
\frac1{g^2(\mu)}=\frac1{g^2(\mu_0)}+
\frac{3}{2(4\pi)^3}\log\frac\mu{\mu_0}.
\]
随着 $\mu$ 增大,$|g|$ 减小并在微扰方程中趋于零;这就是题目要求的渐近自由性质.它由顶点与波函数极点的组合得到,不能用第4套的未计外腿重整化的顶点对数直接替代.
\end{solution}

\originalsection{原题3}
\begin{exercise}
渐近对称性,原题40分.MIT 题单仅引用 Peskin--Schroeder 12.3,并允许在中间步骤使用该书10.2节的结果.
\end{exercise}
\begin{solution}
本次没有取得可再分发的商书完整题面.官方解包含两个标量的四次耦合、固定点及流图,但不能将这些答案反向拼成商书原题的子问.保留原题号与书本定位;完整题面及相应逐子问中文解答为来源缺口,未计入本次可获 MIT 原题的已完成数量.
\end{solution}
